% META: {"id": "TSPE-CONTIN-CR-011", "chapitre": "TSPE-CONTINUITE", "type_objet": "cours", "capacites_codes": ["C2"], "sources_inspiration": [], "mode_creation": "ex_nihilo", "status": "generated"}

\section{Suites définies par récurrence et continuité}

% ============================================================
% STRATE 1 — Theoreme du point fixe
% ============================================================

\theoreme{
Soit $f$ une fonction continue sur un intervalle $I$, telle que $f(I) \subset I$ (c'est-a-dire $f(x) \in I$ pour tout $x \in I$). Soit $(u_n)$ la suite definie par $u_0 \in I$ et $u_{n+1} = f(u_n)$.

Si $(u_n)$ converge vers un reel $\ell \in I$, alors $\ell$ verifie $f(\ell) = \ell$ : $\ell$ est un \textbf{point fixe} de $f$.
}

\demonstration{
Supposons $u_n \to \ell$. Comme $f$ est continue en $\ell$, on a $f(u_n) \to f(\ell)$.

Or $f(u_n) = u_{n+1}$ pour tout $n$, et $u_{n+1} \to \ell$ (meme limite que $(u_n)$, decalee d'un rang).

Par unicite de la limite, $f(\ell) = \ell$.
}

\margeAppui{
Ce theoreme ne dit pas que $(u_n)$ converge : il dit seulement que \emph{si} elle converge, sa limite est necessairement un point fixe de $f$. Il faut etablir la convergence par un autre moyen (monotonie et bornitude, par exemple).
}

\exemple{
Soit $f(x) = \sqrt{2x + 3}$ sur $I = [0, 5]$ et $(u_n)$ definie par $u_0 = 0$, $u_{n+1} = f(u_n)$.

On admet (etude possible par recurrence, cf. chapitre Suites) que $(u_n)$ est croissante et majoree par $3$, donc convergente vers un reel $\ell \in [0, 3]$.

Comme $f$ est continue, $\ell$ verifie $f(\ell) = \ell$, soit $\sqrt{2\ell+3} = \ell$. En elevant au carre (les deux membres sont positifs) : $2\ell + 3 = \ell^2$, soit $\ell^2 - 2\ell - 3 = 0$, c'est-a-dire $(\ell-3)(\ell+1) = 0$. Comme $\ell \geq 0$, on obtient $\ell = 3$.

% BEGIN-VERIFY
% from sympy import symbols, sqrt, solve, Rational
% x = symbols('x')
% sols = solve(sqrt(2*x+3) - x, x)
% assert sols == [3]
% # verification directe de l'equation du second degre
% l = symbols('l')
% assert solve(l**2 - 2*l - 3, l) == [-1, 3]
% END-VERIFY
}

\subsection{Méthode générale}

Pour etudier une suite $(u_n)$ definie par $u_{n+1} = f(u_n)$ avec $f$ continue :
\begin{enumerate}
  \item Montrer que l'intervalle $I$ est stable par $f$ (i.e. $f(I) \subset I$), en general par recurrence.
  \item Etudier la monotonie de $(u_n)$ (souvent par etude du signe de $f(x)-x$, ou par recurrence en comparant $u_{n+1}$ et $u_n$).
  \item Si $(u_n)$ est monotone et bornee, conclure a la convergence par le theoreme de la limite monotone.
  \item Determiner la limite $\ell$ comme solution de $f(\ell) = \ell$ dans $I$ (continuite de $f$).
\end{enumerate}

% ============================================================
% STRATE 2 — Erreurs frequentes
% ============================================================

\erreurFrequente{
\textbf{Chercher le point fixe avant d'avoir prouve la convergence.} L'equation $f(\ell)=\ell$ peut avoir plusieurs solutions dans $I$ ; sans preuve prealable de la convergence de $(u_n)$, on ne peut pas savoir laquelle est la bonne limite (ni meme si la suite converge).
}

\erreurFrequente{
\textbf{Oublier la continuite de $f$.} Le passage a la limite dans $u_{n+1}=f(u_n)$ pour obtenir $\ell=f(\ell)$ n'est valable que si $f$ est continue au point $\ell$. Cette hypothese doit toujours etre verifiee et citee.
}
